\section{Excess Entropy and the \texorpdfstring{$C \times E$}{C x E} Score}\label{app:excess-entropy}

This appendix develops the excess entropy $E$, an alternative summary of the $a_k$ curve that is theoretically interesting (connecting to the excess entropy of computational mechanics \citep{crutchfield2003regularities} and to the total correlation) but empirically inferior to $D_{Ca_n}$ on the external benchmarks of Section~\ref{sec:tevet}.
We include it here for theoretical completeness and to explain the empirical comparisons that motivate our preference for $D_{Ca_n}$.

\subsection{The Excess Entropy}\label{app:excess-entropy-def}

Neither the running average $\bar{a}_n = \frac{1}{n}\sum_k a_k$ nor the total correlation $\TC_n$ is a satisfactory single scalar.
The running average converges to $a_\infty$ (the uninteresting residual floor), losing sensitivity to the learnable structure above it.
The total correlation grows linearly in $n$ and never saturates, because each new term contributes approximately $-\log_2 \theta(r \mid p) - a_\infty > 0$ for large $k$.

The solution is to isolate the signal above the noise floor.
Following the concept of excess entropy from computational mechanics \citep{crutchfield2003regularities}, define the \textbf{excess entropy}:\footnote{The excess entropy $E = \sum_k (a_k - a_\infty)$ is defined as a discrete sum over integer response indices, since ``conditioning on 3.5 responses'' has no direct meaning.
However, a continuous extension exists: consider conditioning on a random subset of size $\alpha n$ for $\alpha \in [0,1]$, yielding a continuous curve $\bar{a}(\alpha)$.
The discrete $\bar{a}_k$ would be samples at $\alpha = k/n$, and $E$ would correspond to integrating this curve (in which case trapezoidal integration would be more faithful than Riemann sums).
We treat $E$ as the discrete sum throughout, but note that the continuous interpretation remains an open question, particularly in connection to the continuous excess entropy literature.
The practical difference between Riemann and trapezoidal sums over $\sim$15 points is small.}
\begin{equation}\label{eq:excess}
    E = \sum_{k=1}^{\infty}\bigl(a_k - a_\infty\bigr)
\end{equation}

Units: bits.
This is the total learnable structural information in the progressive conditioning process, above the irreducible residual noise.
It converges whenever $\theta$'s surprise at new responses eventually stabilizes (since $a_k \to a_\infty$ and the excess $e_k = a_k - a_\infty$ decays to zero).
The ``structure'' captured by $E$ is not limited to discrete mode identity: it includes any inter-response regularity that $\theta$ can learn from context, such as shared stylistic patterns, vocabulary, topic constraints, and formatting conventions.

In practice, $a_\infty$ is unknown.
We estimate it as $a_n$ (the last observed value) and compute:
\begin{equation}\label{eq:excess-hat}
    \hat{E}_n = \sum_{k=1}^{n}(a_k - a_n)
\end{equation}

This underestimates $E$ (since $a_n \geq a_\infty$), but the bias decreases with $n$.

\paragraph{Parametric extrapolation of $a_\infty$.} Rather than using $a_n$ directly, one can fit a parametric model to the observed curve and extrapolate $a_\infty$.
The $a_k$ curve is theoretically expected to be sigmoidal (with concave-up/exponential decay as the degenerate case when the inflection point $k_0 < 1$).
This motivates fitting a four-parameter sigmoid:
\begin{equation}\label{eq:sigmoid-fit}
    a_k = a_\infty + \frac{\alpha}{1 + e^{\beta(k - k_0)}}
\end{equation}
where $a_\infty$ is the asymptotic floor, $\alpha$ is the total drop from $a_1$ to $a_\infty$, $\beta > 0$ controls the steepness of the transition, and $k_0$ is the inflection point.
The fit yields $a_\infty$ without requiring $n$ to be large enough for convergence, and the inflection point $k_0$ provides a principled mode-count estimate when $k_0 > 1$ (Appendix~\ref{app:k0-derivation}).
Section~\ref{sec:mode-count-scaling} provides empirical evidence that the sigmoid-extrapolated $E_{\mathrm{fit}}$ recovers the expected monotonic relationship between mode count and excess entropy, while the raw $\hat{E}_n$ estimator does not.

\paragraph{Relationship to total correlation.} The excess entropy and total correlation are complementary decompositions of the same underlying structure.
Working in expectation over i.i.d.\ draws from $\pi$, let $\bar{a}_1 = \E[-\log_2 \theta(r \mid p)]$ denote the expected unconditional cross-entropy, $\bar{a}_k = \E[a_k]$ the expected conditional cross-entropy at position $k$, and $e_k = \bar{a}_k - a_\infty$ the excess at step $k$.
The per-step mutual information is $I_k = \bar{a}_1 - \bar{a}_k = (\bar{a}_1 - a_\infty) - e_k$.
Summing:
\begin{equation}\label{eq:tc-decomp}
    \TC_n = \sum_{k=1}^{n} I_k = n(\bar{a}_1 - a_\infty) - E_n
\end{equation}

For large $n$ where $E_n \to E$, the total correlation grows linearly with slope $(\bar{a}_1 - a_\infty)$: $\TC_n \approx n(\bar{a}_1 - a_\infty) - E$.
Figure~\ref{fig:decomposition} illustrates.
The three quantities have clean interpretations, all in bits: $(\bar{a}_1 - a_\infty)$ is the per-response redundancy once $\theta$ has fully learned the available structure; $E$ is the total finite structural information, consumed during $\theta$'s transient learning phase; $\TC_n$ is the cumulative redundancy, which grows without bound.

\begin{figure}[tb]
\centering
\begin{tikzpicture}
\begin{axis}[
    width=0.88\textwidth,
    height=6.5cm,
    xlabel={$k$ (response index)},
    ylabel={bits},
    xmin=0.5, xmax=12.5,
    ymin=0, ymax=87,
    xtick={1,2,3,4,5,6,7,8,9,10,11,12},
    ytick={0, 24, 48, 72},
    legend style={font=\small, at={(0.97,0.55)}, anchor=east},
    clip=false,
]

% Define h-bar line
\addplot[name path=hbar, forget plot, draw=none]
    coordinates {(0.5,72) (12.5,72)};

% Define a_infinity line
\addplot[name path=ainf, forget plot, draw=none]
    coordinates {(0.5,24) (12.5,24)};

% Define a_k curve (theoretical, in expectation)
\addplot[name path=ak, forget plot, draw=none]
    coordinates {
        (0.5,72) (1,71) (2,69) (3,66) (4,61)
        (5,53) (6,42) (7,33) (8,28) (9,25)
        (10,24.5) (11,24) (12,24) (12.5,24)
    };

% Fill TC area (between h-bar and a_k) — red
\addplot[red!25, forget plot] fill between[of=hbar and ak];

% Fill E area (between a_k and a_infinity) — blue
\addplot[blue!25, forget plot] fill between[of=ak and ainf];

% Draw the visible lines on top
\addplot[thick, dashed, black]
    coordinates {(0.5,72) (12.5,72)};
\addlegendentry{$\bar{a}_1$ (unconditional)}

\addplot[thick, dotted, black]
    coordinates {(0.5,24) (12.5,24)};
\addlegendentry{$a_\infty$ (floor)}

\addplot[thick, black, mark=*, mark size=1.5pt]
    coordinates {
        (1,71) (2,69) (3,66) (4,61)
        (5,53) (6,42) (7,33) (8,28) (9,25)
        (10,24.5) (11,24) (12,24)
    };
\addlegendentry{$\bar{a}_k$ curve}

% Labels for the colored regions
\node[font=\small] at (axis cs:10,55) {$\TC_n = \sum_k I_k$};
\node[font=\small] at (axis cs:4,36) {$E = \sum_k e_k$};

% Annotation arrows showing I_k and e_k for a single step (k=6, a_6=42)
\draw[<->, thick, gray] (axis cs:6,42) -- node[right, font=\footnotesize] {$e_k$} (axis cs:6,24);
\draw[<->, thick, gray] (axis cs:6,72) -- node[right, font=\footnotesize] {$I_k$} (axis cs:6,42);

\end{axis}
\end{tikzpicture}
\caption{Decomposition of the gap between unconditional surprise $\bar{a}_1$ and the asymptotic floor $a_\infty$ at each step $k$.
The per-step gap $\bar{a}_1 - a_\infty$ splits into mutual information $I_k = \bar{a}_1 - \bar{a}_k$ (\textcolor{red!70}{red}, above the curve) and excess $e_k = \bar{a}_k - a_\infty$ (\textcolor{blue!70}{blue}, below the curve).
Summing across $k$: the red area is the total correlation $\TC_n$; the blue area is the excess entropy $E$.
As $k$ grows, $e_k \to 0$ and each step contributes the full $\bar{a}_1 - a_\infty$ to $\TC_n$, so $\TC_n$ grows linearly while $E$ converges.
All quantities are in bits.}
\label{fig:decomposition}
\end{figure}

\subsection{Per-Byte Excess Entropy Rate}\label{app:per-byte-rate}

The excess entropy $E = \sum_k (a_k - a_\infty)$ is in bits.
For a length-normalized variant, we define the \textbf{per-byte excess entropy rate} by normalizing each response's surprise by its byte count \emph{before} averaging across permutations.
For a single permutation $\sigma$, the per-byte conditional surprise at position $k$ is $a_k^\sigma / \|r_{\sigma(k)}\|$ (bits/byte).
Averaging over permutations gives
\begin{equation}
    \hat{a}_k^{\mathrm{rate}} = \E_\sigma\!\left[\frac{a_k^\sigma}{\|r_{\sigma(k)}\|}\right]
\end{equation}
and the per-byte floor is estimated analogously from the last position: $\hat{a}_\infty^{\mathrm{rate}} = \hat{a}_n^{\mathrm{rate}}$.
The per-byte excess entropy rate is
\begin{equation}\label{eq:per-byte-E}
    E_{\mathrm{rate}} = \sum_{k=1}^{n}\bigl(\hat{a}_k^{\mathrm{rate}} - \hat{a}_\infty^{\mathrm{rate}}\bigr).
\end{equation}
This treats each response equally regardless of length.
Note that $E_{\mathrm{rate}} \neq E / \bar{B}$ in general: because total surprise $a_k$ and byte count $\|r_k\|$ are positively correlated, $E / \bar{B}$ overweights long responses, while the per-response normalization in $E_{\mathrm{rate}}$ avoids this.

\subsection{The \texorpdfstring{$C \times E$}{C x E} Score}

Combining $E$ with the coherence term $C$ (Section~\ref{sec:coherence}) yields scalar scores
\begin{equation}\label{eq:D-total}
    C \times E \quad \text{(bits)} \qquad \text{or} \qquad C \times E_{\mathrm{rate}} \quad \text{(bits/byte)}.
\end{equation}
The first weights total structural information by per-byte plausibility; the second normalizes per byte so long and short responses are treated equally.

\subsection{Why Not Weight Excess Entropy Inside the Sum?}\label{app:why-not-weight}

A natural alternative to reporting $E$ and $C$ separately is to weight the terms within $E$ by coherence, computing $E_w = \sum_k w_k \cdot (a_k - a_\infty)$ where $w_k = 2^{-h_\theta(r_k \mid p)}$.
This would suppress contributions from incoherent modes directly at the point of measurement.
We considered and rejected this approach for three reasons.
First, it breaks the information-theoretic interpretation: $E = \sum_k e_k$ has a clean meaning as the total structural information above the noise floor, connected to Crutchfield and Feldman's excess entropy, and inserting weights makes $E_w$ a hybrid quantity that is not the excess entropy of any well-defined process.
Second, the weights are not independent of what they weight: both $w_k$ and $e_k$ are derived from the same response scored by the same model, and an incoherent response from a novel mode has both high $e_k$ and low $w_k$, which partially cancel in ways that lack a clean characterization.
Third, it makes the noise floor $a_\infty$ ambiguous: unclear whether it should also be weighted, and if so by what.

\subsection{Limitations of $C \times E$}

$C \times E$ inherits two limitations from $E$.
First, $E$ measures learnable \emph{structure} (recurring patterns that $\theta$ can extract from seeing multiple responses) rather than continuous spread.
A policy that draws from a broad, roughly continuous distribution of coherent outputs has $a_k \approx a_1$ for all $k$, yielding $E \approx 0$ despite maximal diversity.
Second, $E$ is near-zero in the few-draws regime: with only 5--10 responses and no repeated modes, $\theta$ finds little learnable structure regardless of diversity, so $C \times E$ carries almost no signal.
These limitations are the reason $D_{Ca_n}$ is preferred as the primary score: $a_\infty$ is the \emph{level} of the floor, which differs meaningfully between diverse and non-diverse response sets even when the curve stays nearly flat.

\paragraph{Empirical inferiority.} On Tevet's diversity-eval benchmark with only 5 responses per set (Section~\ref{sec:tevet}), $C \times a_n$ achieves ROC AUC of \tevetMcDivNugPromptGenCxAnAUC{} on McDiv\_nuggets prompt\_gen while the $D_{\mathrm{fit}} = C \times E_{\mathrm{fit}}$ variant cannot be computed at all (the four-parameter sigmoid fails to converge on only five points) and the discretized $D_{\mathrm{disc}}$ variant is anti-correlated with diversity (AUC = \tevetMcDivNugPromptGenDdiscAUC{}, well below chance).
Figure~\ref{fig:roc-curves} shows the ROC curves.

\begin{figure}[h]
\centering
\includegraphics[width=\textwidth]{tevet_validation/c_ainf_analysis_v3/roc_curves_mcdiv_nuggets.png}
\caption{ROC curves for McDiv\_nuggets binary classification (Qwen2.5-3B, 50 permutations).
Five per-byte metrics are overlaid per panel: $C \times a_n$ (red, solid), $a_n$ alone (orange, dashed), $D_{\mathrm{fit}} = C \times E_{\mathrm{fit}}$ (blue, solid), $D_{\mathrm{disc}} = C \times \hat{E}_n$ (cyan, dashed), and $a_1$ (gray, dotted). Line style is redundant with color so the encoding remains legible under common color-vision deficiencies.
$C \times a_n$ dominates the other four scores on every task; the $E$-based scores ($D_{\mathrm{fit}}, D_{\mathrm{disc}}$) hug or fall below the chance diagonal in this 5-response regime.}
\label{fig:roc-curves}
\end{figure}

\paragraph{Why $E$ fails in this regime.} With only 5 responses and no repeated modes, $\theta$ finds little learnable structure regardless of diversity.
The $a_k$ curve stays approximately flat, yielding $E \approx 0$ for both diverse and non-diverse sets.
The $E$-based scores thus carry almost no signal.
The asymptotic floor $a_\infty$, by contrast, differs substantially between the two conditions: diverse responses remain surprising after conditioning ($a_\infty$ high), while paraphrases become predictable ($a_\infty$ low).
This is the fundamental reason $D_{Ca_n}$ is preferred.

\subsection{Coherence Normalization Strategies}\label{app:c-normalization}

Two strategies are available when coherence variation dominates diversity variation:
\begin{enumerate}[nosep]
    \item \textbf{Normalization}: Replace $C$ with $C / C_{\mathrm{base}}$, where $C_{\mathrm{base}}$ is the coherence of $\theta$'s own samples given the same prompt.
The ratio lives in $(0, 1]$ and equals 1 when $\pi$ is as coherent as $\theta$, removing the baseline suppression.
This has a clean interpretation as ``what fraction of baseline coherence does $\pi$ retain?''
    \item \textbf{Fractional exponent}: Replace $C$ with $C^\alpha$ for $\alpha \in (0, 1)$, which compresses the coherence penalty.
This softens the dominance of $C$ but introduces a free hyperparameter without information-theoretic motivation.
\end{enumerate}
We recommend normalization when a single-scalar ranking is the primary use case, and the unmodified $C$ when the goal is diagnostic.
We did not use either strategy in our experiments.

\subsection{Mode Count Scaling: Excess-Entropy Metrics}

In the synthetic mode count experiments (Section~\ref{sec:mode-count-scaling}), the sigmoid-extrapolated $E_{\mathrm{fit}}$ is monotonic in mode count for Qwen2.5-3B but has wide confidence intervals at high $m$, while the raw $\hat{E}_n$ estimator is non-monotonic.
Table~\ref{tab:mode-count} and Figure~\ref{fig:metrics-vs-m} show the details.

\begin{table}[h]
\centering
\caption{Mode count scaling metrics for Qwen2.5-3B ($n=\modeCountNResponses$, \modeCountNDraws{} draws). $E_{\mathrm{fit}}$ is the sigmoid-extrapolated excess entropy; $a_n$ is the mean floor at $k=\modeCountNResponses$; $\sigma_\ell$ is per-byte coherence spread; $k_0$ is the sigmoid inflection point.
Confidence intervals are 95\% bootstrap.}
\label{tab:mode-count}
\small
% Generated by: scripts/generate_mode_count_table.py
% Data source: results/mode_count/qwen2.5-3b_1k_draws.json
%              results/mode_count/qwen2.5-3b_fits.json
\begin{tabular}{@{}r rrrr@{}}
\toprule
$m$ & $E_{\mathrm{fit}}$ (bits) & $a_n$ (bits) & $\sigma_\ell$ (bits/byte) & $k_0$ \\
\midrule
1  & $296 \pm 11$ & 9.3 & 0.1 & $-10.0$ \\
2  & $568 \pm 24$ & 16.9 & 0.2 & $-10.0$ \\
3  & $704 \pm 40$ & 26.7 & 0.2 & $-10.0$ \\
4  & $799 \pm 67$ & 37.4 & 0.3 & $-10.0$ \\
5  & $855 \pm 113$ & 48.0 & 0.3 & $-10.0$ \\
6  & $940 \pm 127$ & 56.1 & 0.3 & $-10.0$ \\
7  & $1069 \pm 194$ & 61.5 & 0.3 & $-10.0$ \\
8  & $1138 \pm 253$ & 66.9 & 0.3 & $-10.0$ \\
9  & $1563 \pm 705$ & 72.8 & 0.3 & $-10.0$ \\
10  & $1807 \pm 1194$ & 77.8 & 0.3 & $-10.0$ \\
\bottomrule
\end{tabular}

\end{table}

\begin{figure}[h]
\centering
\includegraphics[width=0.6\textwidth]{mode_count/qwen2.5-3b/metrics_vs_m.png}
\caption{Summary metrics vs.\ mode count $m$ (Qwen2.5-3B, $n=\modeCountNResponses$, \modeCountNDraws{} draws). $E_{\mathrm{fit}}$ (sigmoid-extrapolated) increases monotonically, while the raw $\hat{E}_n$ does not.}
\label{fig:metrics-vs-m}
\end{figure}

\begin{figure}[h]
\centering
\includegraphics[width=0.6\textwidth]{mode_count/qwen2.5-3b/fit_params_vs_m.png}
\caption{Sigmoid fit parameters vs.\ mode count $m$ (Qwen2.5-3B).
The inflection point $k_0$ remains at the lower bound ($-10$) across all $m$, indicating pure exponential decay without an initial plateau (see Section~\ref{sec:cross-mode-learning}).}
\label{fig:fit-params}
\end{figure}

% Generated by scripts/dataset_confound_analysis.py
% Per-byte a_1 for r_0 of the illustrative high- and low-diversity samples.
% high sample: test.sa-sb.sa::00698
% low sample:  test.sa-sb.sa-b::00279
\newcommand{\confoundHighRoAone}{0.783}
\newcommand{\confoundLowRoAone}{1.777}
\newcommand{\confoundHighSampleId}{test.sa-sb.sa::00698}
\newcommand{\confoundLowSampleId}{test.sa-sb.sa-b::00279}

